\subsection{Exterior support of the Jacobian component}
\label{sec:schur-support}

The obstruction to recombining two projective Schur components has a
simple linear-algebraic description.  A secant or tangent vector of
\(\Gr(4,W)\) uses at most eight directions in \(W\).  The
Jacobian component of a smooth quartic uses all ten.  We prove a
sharper statement which also applies to singular quartics.

Fix a volume \(\epsilon=x_1\wedge x_2\wedge x_3\wedge x_4\)
and suppress it when writing quartics.  Define the equivariant map
\[
 j:\det V\otimes\Sym^4V\longrightarrow\bigwedge^4\Sym^2V
\]
by polarization of
\begin{equation}\label{eq:polarized-schur-embedding}
 j(\epsilon\otimes\ell^4)
       =(\ell x_1)\wedge(\ell x_2)\wedge(\ell x_3)\wedge(\ell x_4).
\end{equation}
The right side changes by the determinant of a change of the
\(x_i\), so this definition does not depend on the chosen volume.
More explicitly,
\begin{equation}\label{eq:schur-polarization-formula}
 j(\epsilon\otimes v_1v_2v_3v_4)
 =\frac1{24}\sum_{\sigma\in\mathfrak S_4}
       (v_{\sigma(1)}x_1)\wedge\cdots\wedge(v_{\sigma(4)}x_4).
\end{equation}
This is nonzero.  Its irreducible source and
\eqref{eq:two-web-components} identify its image with \(E_{35}\).
For the Jacobian map \(C\), evaluation at \(\ell=x_1\) gives
\(Cj(x_1^4)=2x_1^4\); hence
\begin{equation}\label{eq:jacobian-schur-normalization}
 Cj=2\,\mathrm{id},\qquad Q=\tfrac12jC.
\end{equation}
These identities also give \(j=C^*/24\) in the Fischer
normalization used below.

For \(q\in W^*=\Sym^2V^*\), regard \(q\) as the symmetric
bilinear functional \(q(uv)\), and let \(\iota_q\) denote
exterior contraction.  A quartic has \emph{essential-variable
space} \(H\subset V\) if \(H\) is the smallest subspace for
which it lies in \(\Sym^4H\).  Equivalently,
\(H^\perp\) is the kernel of its first differentiation map
\(V^*\to\Sym^3V\).  This is the first-catalecticant characterization of essential variables
\cite[Proposition~1]{Carlini}; it also follows by choosing
coordinates along the kernel and using characteristic zero.
This subspace of \(V\) is distinct from the exterior-support
subspace of \(W=\Sym^2V\) defined below.

\subsubsection{Two representation-theoretic lemmas}

We keep the determinant line \(D=\det V\) in the notation.
The first lemma concerns a polynomial representation with its actual
weight decomposition, not just its associated graded vector space.

\begin{lemma}[Shear descent in a fixed weight]
\label{lem:shear-descent}
Let \(V=A\oplus B\), with \(B\ne0\), and let \(m\ge1\).
Suppose that \(N\subset D\otimes\Sym^mV\) is invariant under
\(\mathrm{id}_A\oplus t\,\mathrm{id}_B\), \(t\in\C^*\),
and under all shears
\[
 u_L(a+b)=a+b+L(a),\qquad L\in\Hom(A,B).
\]
If \(N\not\subset D\otimes\Sym^mB\), then
\[
 N\cap\bigl(D\otimes A\otimes\Sym^{m-1}B\bigr)\ne0.
\]
More precisely, a nonzero vector in the \(A\)-degree-\(k\)
weight space, \(k>0\), has a nonzero image in the degree-one
space under a product of \(k-1\) infinitesimal rank-one shears.
The same fixed nonzero vector of \(B\) can be used in every shear.
\end{lemma}
\begin{proof}
Put \(s=\dim B\).  The displayed torus acts on
\(D\otimes\Sym^kA\otimes\Sym^{m-k}B\) with weight
\(s+m-k\).  These weights are distinct.  Since a complex torus
acts semisimply, its invariant subspace \(N\) is the direct sum
of its intersections with these weight spaces.  This justifies
projecting a vector of \(N\) to one degree before differentiating;
no cancellation between degrees is possible.

Choose a basis \(a_1,\ldots,a_r\) of \(A\), a nonzero
\(b\in B\), and a basis vector \(\epsilon\) of \(D\).
Write a nonzero degree-\(k\) vector as \(\epsilon\otimes F\),
where
\[
 F=\sum_{|\beta|=k}a^\beta f_\beta,\qquad
                 f_\beta\in\Sym^{m-k}B.
\]
Choose \(\beta\) with \(f_\beta\ne0\), and \(i\) with
\(\beta_i>0\).  Set \(\alpha=\beta-e_i\).  In
\(\partial_A^\alpha F\), the coefficient of \(a_i\) is
exactly \(\beta!f_\beta\).  Indeed a monomial that contributes
to this coefficient must have exponent \(\alpha+e_i=\beta\).
This coefficient is nonzero in characteristic zero, so the derivative
is a nonzero element of \(A\otimes\Sym^{m-k}B\).

The shear sending \(a_j\) to \(a_j+tb\) and fixing the other
basis vectors has determinant one.  Its infinitesimal action on
\(D\otimes\Sym^mV\) is consequently
\(1_D\otimes b\partial_{a_j}\), with no extra action on
\(D\).  Invariance under the shear implies invariance under its
infinitesimal action: the orbit is a polynomial in \(t\), and
all its coefficients belong to the invariant linear subspace.
These differential operators commute, because every
\(\partial_{a_j}\) kills \(b\).  Their product with
multiplicities \(\alpha_j\) sends our vector to
\[
 \epsilon\otimes b^{k-1}\partial_A^\alpha F.
\]
This vector belongs to \(N\).  It is nonzero, since multiplication
by \(b^{k-1}\) is injective in the polynomial algebra \(\Sym B\).
Its exact bidegree is \((1,m-1)\); thus its image in the corresponding
weight-graded piece is also nonzero.  The proof uses neither a special
choice of \(b\) nor a restriction on \(r\) or \(s\).  When
\(k=1\) the product of shears is the identity.  This proves all
assertions, including the low-dimensional cases.
\end{proof}

For the next lemma write
\([a,b]=\Sym^a U\otimes\Sym^b U'\), where \(U,U'\)
are the two standard two-dimensional representations of
\(SL_2\times SL_2\).  Only even weights occur below, so all
representations descend through its double cover of \(SO_4\).

\begin{lemma}[The exterior-cube character]
\label{lem:orthogonal-exterior-cube}
Let \(q\) be nondegenerate on the four-dimensional space \(V\),
and put \(W=\Sym^2V\).  Under \(SL_2\times SL_2\to SO(V,q)\),
choose \(V=[1,1]\).  Then
\begin{equation}\label{eq:explicit-so4-exterior-cube}
\begin{split}
 \bigwedge^3 W\simeq{}&[6,0]\oplus[0,6]\oplus[4,4]
       \oplus2[4,2]\oplus2[2,4]\\
       &\oplus[2,2]\oplus2[2,0]\oplus2[0,2].
\end{split}
\end{equation}
In particular \((\bigwedge^3W)^{SO(V,q)}=0\), and
\(\Hom_{O(V,q)}(D,\bigwedge^3W)=0\).
\end{lemma}
\begin{proof}
Put \(A_2=\Sym^2U\) and \(B_2=\Sym^2U'\).
The symmetric Cauchy formula gives
\[
 W=(A_2\otimes B_2)\oplus\C\rho,
                 \qquad \rho=q^{-1},
\]
where \(\rho\) is invariant.  Write \(T=A_2\otimes B_2\).
For any two three-dimensional vector spaces, antisymmetrization in
\((A_2\otimes B_2)^{\otimes r}\), for \(r=2,3\), gives
\begin{align*}
 \bigwedge^2T&=(\Sym^2A_2\otimes\bigwedge^2B_2)
                 \oplus(\bigwedge^2A_2\otimes\Sym^2B_2),\\
 \bigwedge^3T&=(\Sym^3A_2\otimes\bigwedge^3B_2)
       \oplus(\mathbb S_{(2,1)}A_2\otimes\mathbb S_{(2,1)}B_2)
       \oplus(\bigwedge^3A_2\otimes\Sym^3B_2).
\end{align*}
For clarity, the reason for the conjugate partitions in these formulas
is that tensoring a symmetric-group representation with the sign
representation conjugates its partition.  Schur--Weyl decomposition
of the two tensor powers therefore gives exactly the displayed
summands.  This is the usual skew Cauchy formula, not an additional
hypothesis about this particular \(q\).

Here are all the required one-factor calculations:
\begin{equation}\label{eq:spin-one-calculations}
\begin{array}{c|c}
\text{functor on }\Sym^2U&\text{restriction to }SL_2\\ \hline
\Sym^2&\Sym^4U\oplus\C\\
\bigwedge^2&\Sym^2U\\
\bigwedge^3&\C\\
\Sym^3&\Sym^6U\oplus\Sym^2U\\
\mathbb S_{(2,1)}&\Sym^4U\oplus\Sym^2U
\end{array}
\end{equation}
The weights of \(\Sym^2U\) are \(2,0,-2\).
Taking unordered pairs or triples of these weights verifies the
symmetric-power rows; taking distinct weights verifies the exterior
rows.  Complete reducibility identifies the resulting characters
with the stated representations.  Alternatively the final row
follows from
\(A_2\otimes\bigwedge^2A_2=
\mathbb S_{(2,1)}A_2\oplus\bigwedge^3A_2\): the tensor product
\(\Sym^2U\otimes\Sym^2U\) is
\(\Sym^4U\oplus\Sym^2U\oplus\C\), and the exterior cube
removes the scalar summand.

Substitution gives
\begin{align*}
 \bigwedge^2T&=[4,2]\oplus[0,2]\oplus[2,4]\oplus[2,0],\\
 \bigwedge^3T&=[6,0]\oplus[2,0]\oplus[4,4]\oplus[4,2]
       \oplus[2,4]\oplus[2,2]\oplus[0,6]\oplus[0,2].
\end{align*}
Since \(\bigwedge^3W=\bigwedge^3T\oplus
(\bigwedge^2T)\wedge\rho\), these equalities prove
\eqref{eq:explicit-so4-exterior-cube}.  The dimensions are
\(84+36=120=\binom{10}{3}\).  There is no summand \([0,0]\).
Finally the determinant character of \(O(V,q)\), namely its
representation on \(D\), restricts to the trivial character of
\(SO(V,q)\).  Any equivariant image of \(D\) would therefore
be a nonzero \(SO(V,q)\)-invariant vector, which the decomposition
excludes.  This tracks the determinant twist without choosing an
extension of each nontrivial \(SO_4\) summand to \(O_4\).
\end{proof}


\begin{lemma}[The contraction kernel]
\label{lem:schur-contraction-kernel}
For \(0\ne q\in\Sym^2V^*\), let
\(\kappa_q:D\otimes\Sym^4V\to\bigwedge^3\Sym^2V\)
be the map \(\epsilon\otimes f\mapsto\iota_qj(\epsilon\otimes f)\).  If \(q\) is singular, then
\begin{equation}\label{eq:singular-contraction-kernel}
 \ker\kappa_q=\det V\otimes\Sym^4(\rad q).
\end{equation}
If \(q\) is nondegenerate, its inverse is a quadratic tensor
\(q^{-1}\in\Sym^2V\), and
\begin{equation}\label{eq:nonsingular-contraction-kernel}
 \ker\kappa_q=\det V\otimes\C(q^{-1})^2.
\end{equation}
\end{lemma}
\begin{proof}
For each fixed \(q\), the map in the statement is explicitly
\[
 \kappa_q:D\otimes\Sym^4V\longrightarrow\bigwedge^3\Sym^2V,
              \qquad \epsilon\otimes f\longmapsto
                     \iota_qj(\epsilon\otimes f).
\]
Naturality of polarization and contraction makes this map equivariant
for the stabilizer of \(q\) in \(GL(V)\).  The source carries
the determinant factor \(D\); the target has its ordinary exterior
representation.  We use that equivariance in both cases.

Suppose first that \(q\) has rank \(r\in\{1,2,3\}\).
Choose \(V=A\oplus B\), where \(B=\rad q\) and \(q|_A\)
is nondegenerate.  Every pure power \(b^4\), \(b\in B\),
is killed by contraction in
\eqref{eq:polarized-schur-embedding}, since \(q(bx)=0\) for
all \(x\in V\).  Pure powers span \(\Sym^4B\), so
\(D\otimes\Sym^4B\subset\ker\kappa_q\).

The stabilizer contains \(O(A,q|_A)\times GL(B)\), the
\(B\)-scaling torus, and all shears in
Lemma~\ref{lem:shear-descent}.  On the degree-one weight space
\[
 D\otimes A\otimes\Sym^3B
\]
the product-group representation is irreducible.  Indeed its untwisted
factors are the standard representation of the full orthogonal group
and the irreducible symmetric cube for \(GL(B)\).  The determinant
factor restricts to the product character \(\det_A\det_B\),
which does not change irreducibility.  For \(r=1\) the standard
factor is one-dimensional; for \(r=2\) its two isotropic weight
lines for \(SO_2\) are exchanged by a reflection in \(O_2\),
so neither is invariant under the full group.  The usual standard
representation is irreducible also for \(r=3\).

The restriction of \(\kappa_q\) to this space is not zero.
Choose an orthogonal basis with \(a=x_1\in A\),
\(q(a,a)=1\), and \(b=x_4\in B\), and set
\(\epsilon=x_1\wedge\cdots\wedge x_4\).
In the coefficient of \(st^3\) in
\(\iota_qj(\epsilon\otimes(sa+tb)^4)\), only the contraction
of the first factor contributes.  Thus
\[
 4\kappa_q(\epsilon\otimes ab^3)
       =(bx_2)\wedge(bx_3)\wedge b^2\ne0.
\]
The three factors are independent because multiplication by \(b\)
is injective on \(V\).  Irreducibility now gives injectivity on
the degree-one space.  If the kernel contained any other positive
\(A\)-degree, Lemma~\ref{lem:shear-descent}, with \(m=4\),
would give a nonzero kernel vector in that degree-one space, a
contradiction.  This proves \eqref{eq:singular-contraction-kernel}
for every singular rank, without a genericity condition on \(q\).

Now let \(q\) be nondegenerate.  After a change of coordinates,
\(q(x_ix_j)=\delta_{ij}\) and
\(\rho=q^{-1}=x_1^2+x_2^2+x_3^2+x_4^2\).
The harmonic decomposition is
\begin{equation}\label{eq:orthogonal-quartic-decomposition}
 D\otimes\Sym^4V=
 (D\otimes\mathcal H_4)\oplus
 (D\otimes\rho\mathcal H_2)\oplus(D\otimes\C\rho^2).
\end{equation}
Here \(\mathcal H_m\) is the kernel of the Laplacian in degree
\(m\).  To check the decomposition directly, the symmetric Cauchy
formula for \(V=U\otimes U'\) in degree four has partitions
\((4),(3,1),(2,2)\), giving \(SO_4\)-types
\([4,4],[2,2],[0,0]\) of dimensions \(25,9,1\).
Multiplication by the invariant \(\rho\) identifies the degree-two
and scalar pieces; differentiation verifies the harmonic description.
The three summands are irreducible and pairwise nonisomorphic for
\(SO_4\).  The determinant factor restricts trivially to this group.

The last summand is killed: as an \(O(V,q)\)-representation,
\(D\otimes\C\rho^2\simeq D\), and
Lemma~\ref{lem:orthogonal-exterior-cube} proves that the target of
\(\kappa_q\) has no such character.  It remains to show that the
other two restrictions are nonzero.  Put
\[
 h_4=x_1^4-6x_1^2x_2^2+x_2^4\in\mathcal H_4,
            \qquad h_2=x_1^2-x_2^2\in\mathcal H_2.
\]
In the basis triple
\(\tau=(x_1x_2)\wedge(x_1x_3)\wedge(x_1x_4)\),
formula~\eqref{eq:schur-polarization-formula} gives
\[
 [\tau]\kappa_q(\epsilon\otimes h_4)=2,
     \qquad [\tau]\kappa_q(\epsilon\otimes\rho h_2)=\frac23.
\]
For the first equality the terms \(x_1^4\) and
\(-6x_1^2x_2^2\) contribute \(1,1\).  For the second the
terms \(x_1^4,x_1^2x_3^2,x_1^2x_4^2\) contribute
\(1,-1/6,-1/6\); all other terms contribute zero.
These are coefficient identities, not numerical rank assumptions.
Each nonzero restriction from an irreducible summand is injective,
and the images of the two nonisomorphic summands cannot intersect.
The kernel is therefore precisely \(D\otimes\C\rho^2\).
Equivariance transports this conclusion back to the original form
\(q\), proving \eqref{eq:nonsingular-contraction-kernel}.
\end{proof}


For \(z\in\bigwedge^4W\), write \(s(z)\) for the dimension
of its smallest exterior support: the smallest subspace \(S\)
with \(z\in\bigwedge^4S\).  In any basis,
\begin{equation}\label{eq:exterior-support-annihilator}
 s(z)=10-\dim\{q\in W^*: \iota_qz=0\}.
\end{equation}
Indeed the vanishing of contraction by a dual basis vector is exactly
the absence of all wedge monomials containing its corresponding
basis vector.  Applying this after a change of basis proves both the
formula and the existence of the smallest support.

\begin{proposition}[Exact support of the Schur embedding]
\label{prop:exact-schur-support}
Let \(0\ne f\in\Sym^4V\) have essential-variable dimension
\(d\).  With the determinant factor understood,
\begin{equation}\label{eq:schur-support-table}
\begin{array}{c|l|c}
 d&\text{condition on }f&s(j(f))\\ \hline
 1&\text{arbitrary}&4\\
 2&\text{arbitrary}&7\\
 3&\text{arbitrary}&9\\
 4&f\text{ not the square of a nondegenerate quadratic}&10\\
 4&f\text{ the square of a nondegenerate quadratic}&9
\end{array}
\end{equation}
In particular, \(s(j(f))\ge9\) if \(d\ge3\), and it is
\(10\) if the projective quartic \(V(f)\) is smooth.
\end{proposition}
\begin{proof}
If \(d<4\) and \(H\) is the essential-variable space, the
contraction kernel lemma says that the annihilating bilinear forms
are exactly those with \(H\subset\rad q\).  They form
\(\Sym^2(V/H)^*\), of dimension \((4-d)(5-d)/2\).
No nondegenerate form can annihilate \(j(f)\), since the square
of a nondegenerate quadratic has four essential variables.
This gives the first three entries of
\eqref{eq:schur-support-table}.

If \(d=4\), no singular form annihilates \(j(f)\).  A
nondegenerate form does so exactly when \(f\) is proportional
to its inverse square.  In this case its projective class is unique:
proportional squares of nonzero quadratics over \(\C\) have
proportional square roots, and inversion is projectively unique.
The annihilator therefore has dimension one in the square case and
zero otherwise.  Formula~\eqref{eq:exterior-support-annihilator}
gives the last two entries.

A quartic with fewer than four essential variables is a cone and is
singular at its vertex.  A square of a quadratic is nonreduced and
has vanishing gradient along that quadratic.  Neither defines a
smooth quartic surface, proving the last assertion.
\end{proof}

\begin{lemma}[The eight-dimensional secant and tangent bound]
\label{lem:secant-tangent-support}
A sum of two decomposable four-vectors in \(\bigwedge^4W\)
has exterior support at most eight.  Every vector in the affine
tangent space \(\bigwedge^3R\wedge W\) to the Pluecker cone
at \(w_R\ne0\) also has exterior support at most eight.
\end{lemma}
\begin{proof}
For a sum, take the sum of the two supporting four-planes.  For a
tangent vector, choose a basis \(r_1,\ldots,r_4\) of \(R\).
Every element of \(\bigwedge^3R\wedge W\) has the form
\[
 \sum_{i=1}^4 r_1\wedge\cdots\wedge r_{i-1}
             \wedge u_i\wedge r_{i+1}\wedge\cdots\wedge r_4
\]
for four vectors \(u_i\in W\), and is supported in
\(R+\langle u_1,\ldots,u_4\rangle\).
\end{proof}

\begin{theorem}[No exceptional recombination above three essential variables]
\label{thm:global-schur-recombination}
Let \(R\in\Gr(4,\Sym^2V)\) have nonzero Jacobian quartic
\(f_R=Cw_R\) with at least three essential variables.  Then
both \(Pw_R,Qw_R\) are nonzero and
\[
 \dim H_{w_R}=2,\qquad
 \mathscr L_R\cap\Gr(4,W)=\{[w_R]\}
\]
as schemes.  In particular,
\begin{equation}\label{eq:global-reconstruction-open-equality}
                  G_4^{\mathrm{rec}}=G_4^\circ.
\end{equation}
The secant, double-point, pure-endpoint and flag-line strata of the
component pencil are all disjoint from \(G_4^\circ\).
\end{theorem}
\begin{proof}
By \eqref{eq:jacobian-schur-normalization} and
Proposition~\ref{prop:exact-schur-support}, the vector \(Qw_R\)
has support at least nine.  It is nonzero and cannot be proportional
to the decomposable vector \(w_R\); hence \(Pw_R\ne0\).
On the component line, the quadratic Pluecker ideal restricts to
\((b-a)H_{w_R}\), where
\(H_{w_R}=\langle\ell_{F,R}\rangle\subset H^0(\Pj^1,\OO(1))\).
If this space has dimension two, its ideal sheaf is \((b-a)\).
If it has dimension one, the intersection is the degree-two divisor
of \((b-a)\ell\), including multiplicities.  If it is zero,
the entire line is contained in the Grassmannian.  These assertions
follow directly from the ideal generated by the Pluecker quadrics;
no classification of an ambient boundary is used.

If the closed pencil has a distinct second decomposable point
\(z=aPw_R+bQw_R\) with \(a\ne b\), then
\(z=a w_R+(b-a)Qw_R\).  Thus \(Qw_R\) is a linear
combination of two decomposable vectors, contradicting
Lemma~\ref{lem:secant-tangent-support}.  This argument includes a
pure endpoint.  If the residual point coincides with \([w_R]\),
then the line is tangent at \([w_R]\): its restricted equations
have zero first derivative there.  Its direction is \(Qw_R\)
modulo \(\C w_R\), so \(Qw_R\in\bigwedge^3R\wedge W\),
giving the same contradiction.  A whole Grassmann line also has that tangent
property.  All alternatives to rank two are excluded.

Smoothness of \(Y_R\) implies four essential variables, so every
\(R\in G_4^\circ\) satisfies this conclusion.  The equality
in \eqref{eq:global-reconstruction-open-equality} holds as an
equality of open subschemes.  Here \(G_4^\circ\) has the
reduced open-subscheme structure inherited from \(\Gr(4,10)\).
The complement is defined by the two-by-two minors of the
two-column coefficient matrix of the residual linear forms.
It is a closed subscheme of finite type over \(\C\) with no
geometric points, hence is empty, regardless of possible nilpotents
in that determinantal presentation.  The statement for all the named
exceptional strata follows at the same time.
\end{proof}

\begin{corollary}[An invariant location of the ambient obstruction]
\label{cor:exceptional-binary-jacobian}
On \(X^\times\), the support of the rank-defect locus
\(\dim H_w\le1\) is contained in the inverse image under the
Jacobian covariant of the subspace variety of binary quartics.  That
variety is defined set-theoretically by rank at most two of the first
catalecticant \(V^*\to\Sym^3V\).  In particular, it produces
no exceptional divisor or ramification locus in the smooth geometric
family, and no pair of non-projectively-equivalent smooth webs with
isomorphic full failure schemes.
\end{corollary}
\begin{proof}
The containment is the contrapositive of
Theorem~\ref{thm:global-schur-recombination}; the catalecticant
characterization is the essential-variable characterization above.
The inverse-scheme assertion follows from the intrinsic readout and
Theorem~\ref{thm:generic-intrinsic-web}.  No equality with the whole
binary-Jacobian inverse image, or classification of its ambient
irreducible components, is asserted.
\end{proof}
